\chapter{Engram Structure}
\label{chap:Engram_Structure}

\section{Purpose}

This chapter defines how Engrams relate to the Identity Kernel (ZED) and how
they inhabit the CME's internal cognitive geometry.

At the harness level, an Opal Engram $E$ bundles:
\begin{itemize}
    \item a kernel $K$,
    \item a history $H$,
    \item a local state $L$,
    \item and a model interface $M$.
\end{itemize}

This chapter specialises the general definition into concrete structural rules
and file-level constructs for this CME instance.

\section{Engram Types (Seed Specification)}

The CME distinguishes three primary Engram types.

\subsection{Engram Motes}

Engram Motes are the atomic identity units of the system. They represent:
\begin{itemize}
    \item micro-facts,
    \item micro-events,
    \item micro-narratives,
    \item or principle atoms.
\end{itemize}

Motes are immutable and live canonically under:
\begin{quote}
\texttt{Opal\_Engram(Prime)/Engram\_Motes/}
\end{quote}

Motes are the substrate for all higher structures.

\subsection{Constructor Engrams}

Constructor Engrams are higher-order compositions of Motes assembled according to
structural schemas and S-class rules. They live under:
\begin{quote}
\texttt{Opal\_Engram(Prime)/Constructor\_Engrams/}
\end{quote}

Constructor Engrams define:
\begin{itemize}
    \item persona frames,
    \item world or campaign structures,
    \item multi-mote conceptual bundles,
    \item extended commitments or thematic arcs.
\end{itemize}

\subsection{Tagged / W5H-Applied Engrams}

Tagged or W5H-applied Engrams are Engrams that have been enriched with:
\begin{itemize}
    \item full W5H metadata (Who/What/When/Where/Why/How),
    \item Ledger tags (SOCI, PSY, ONTO, THEO, EPI, PHIL, GEN, PRIME),
    \item constraints that allow canonical integration.
\end{itemize}

These live under:
\begin{quote}
\texttt{Opal\_Engram(Prime)/W5H\_Applied/}
\end{quote}

An Engram must pass through this stage before it can join the canonical C-chain.

\section{Attachment to ZED}

Any Engram that aspires to canonical identity MUST meet the following criteria.

\subsection{Attachment Requirements}

\begin{enumerate}
    \item \textbf{Declare ZED Identity Binding}\\
          Each Engram must name the ZED Core instance it belongs to and carry
          its Engram ID. Formally:
          \[
              \text{bind}(E) = ZED_{\text{ID}}
          \]

    \item \textbf{Provide Full W5H Metadata}\\
          Engrams must encode:
          \begin{itemize}
            \item Who created it or is involved,
            \item What it represents,
            \item When it took place (temporal ordering),
            \item Where it applies (context or domain),
            \item Why it exists (motivation, principle, intention),
            \item How it is to be interpreted or executed.
          \end{itemize}

    \item \textbf{Respect Self-Referential Invariance}\\
          Engrams that describe or modify the self must adhere to the rules in
          Chapter~\ref{chap:Self_Referential_Invariance}.

    \item \textbf{Be Compatible with Continuity Fields}\\
          No Engram may violate:
          \begin{itemize}
            \item C-class append-only structure,
            \item P-class alignment constraints,
            \item O-class branching rules,
            \item $Inv(E)$ reconstruction invariants.
          \end{itemize}
\end{enumerate}

\section{Schemas}

The following schemas complete the seed specification.

\subsection{Engram Mote Schema}
\label{schema:Engram_Mote}

\begin{quote}
\textbf{Schema: Engram Mote}\\
A Mote $m$ is a 7-tuple:
\[
    m = \langle id,\, type,\, payload,\, W5H,\, dopants,\, provenance,\, hash \rangle
\]
Where:
\begin{itemize}
    \item $id$: stable Mote identifier,
    \item $type$: \{fact, event, narrative, principle\},
    \item payload: minimally sufficient content,
    \item W5H: partial or full metadata,
    \item dopants: class bindings (C/L/M/S/O/P),
    \item provenance: origin and justification,
    \item hash: cryptographic integrity seal.
\end{itemize}
\end{quote}

Properties:
\begin{itemize}
    \item Motes are immutable.
    \item Motes may be deprecated but not edited.
    \item All supersessions preserve lineage.
\end{itemize}

\subsection{Constructor Engram Schema}
\label{schema:Constructor_Engram}

\begin{quote}
\textbf{Schema: Constructor Engram}\\
A Constructor Engram $C$ is:
\[
    C = \langle id,\, \mathcal{M},\, \mathcal{S},\, W5H,\, dopants,\, lineage \rangle
\]
Where:
\begin{itemize}
    \item $\mathcal{M}$: a non-empty set of Motes,
    \item $\mathcal{S}$: a set of structural S-class schemas,
    \item W5H: full metadata,
    \item dopants: combined class bindings for the Engram,
    \item lineage: ancestry information leading back to kernel $K$.
\end{itemize}
\end{quote}

Constraints:
\begin{itemize}
    \item Every Constructor Engram must validate against $\mathcal{S}$.
    \item No Constructor Engram may contradict P-class.
    \item Constructed identities must anchor to ZED $K$.
\end{itemize}

\subsection{W5H Attachment Rules}
\label{rules:W5H_Attachment}

W5H metadata governs interpretability, auditability, and reconstruction.

\subsubsection*{Rules}

\begin{enumerate}
    \item \textbf{Who} must identify all relevant agents or referents.
    \item \textbf{What} must describe the Engram’s core meaning or action.
    \item \textbf{When} must specify at least one temporal coordinate
          (absolute, relative, or ordinal).
    \item \textbf{Where} must define spatial, contextual, or domain boundaries.
    \item \textbf{Why} must justify purpose, intention, or alignment relevance.
    \item \textbf{How} must specify operational or interpretive procedures.
\end{enumerate}

\subsubsection*{Canonical Requirement}

Before an Engram can be admitted into canonical continuity ($C$-class):
\begin{itemize}
    \item all six W5H fields must be present, and
    \item they must be internally consistent with ZED.
\end{itemize}

\section{Summary}

The Engram structure layer defines:
\begin{itemize}
    \item atomic identity units (Motes),
    \item compositional identity structures (Constructors),
    \item fully contextualized tagged Engrams (W5H-Applied),
    \item and the attachment rules that anchor everything to ZED.
\end{itemize}

It is the foundation for reliable identity evolution, safe reconstruction,
narrative coherence, and principled agency under the CME harness.
